\begin{abstract}
    Animal re-identification (ReID) faces critical challenges due to viewpoint variations, particularly in Aerial-Ground (AG-ReID) settings where models must match individuals across drastic elevation changes.
    However, existing datasets lack the precise angular annotations required to systematically analyze these geometric variations.

    To address this, we introduce the \textbf{M}ulti-view \textbf{O}riented \textbf{O}bservation (MOO) dataset, a large-scale synthetic AG-ReID dataset of $1,000$ cattle individuals captured from $128$ uniformly sampled viewpoints ($128,000$ annotated images).

    Using this controlled dataset, we quantify the influence of elevation and identify a critical elevation threshold, above which models generalize significantly better to unseen views.

    Finally, we validate the transferability to real-world applications in both zero-shot and supervised settings, demonstrating performance gains across four real-world cattle datasets and confirming that synthetic geometric priors effectively bridge the domain gap.
    Collectively, this dataset and analysis lay the foundation for future model development in cross-view animal ReID.

    MOO is publicly available at \url{https://github.com/TurtleSmoke/MOO}.
\end{abstract}